\documentclass[11pt,a4paper,openany]{book}
\usepackage[margin=1in]{geometry}
\usepackage{amsmath,amssymb,amsthm,mathtools}
\usepackage{enumitem}
\usepackage{booktabs,array,longtable}
\usepackage{tikz}
\usepackage{microtype}
\microtypesetup{expansion=false}
\usepackage[colorlinks=true,linkcolor=black!70!blue,citecolor=black!70!blue,urlcolor=black!70!blue]{hyperref}
\setlist{itemsep=2pt,topsep=4pt}
\setlength{\parindent}{0pt}
\setlength{\parskip}{5pt plus 1pt}
\renewcommand{\arraystretch}{1.25}
\setcounter{tocdepth}{2}

% ---------- theorem environments (classic amsthm) ----------
\theoremstyle{plain}
\newtheorem{theorem}{Theorem}[chapter]
\newtheorem{lemma}[theorem]{Lemma}
\newtheorem{proposition}[theorem]{Proposition}
\newtheorem{corollary}[theorem]{Corollary}
\newtheorem{claim}[theorem]{Claim}
\theoremstyle{definition}
\newtheorem{definition}[theorem]{Definition}
\newtheorem{example}[theorem]{Example}
\newtheorem{exercise}[theorem]{Exercise}
\newtheorem{axiom}{Axiom}
\theoremstyle{remark}
\newtheorem{remark}[theorem]{Remark}
\newtheorem*{intuition}{Intuition}
\newtheorem*{warning}{Warning}
\newtheorem*{notation}{Notation}

% ---------- macros ----------
\newcommand{\N}{\mathbb{N}}
\newcommand{\Z}{\mathbb{Z}}
\newcommand{\Q}{\mathbb{Q}}
\newcommand{\R}{\mathbb{R}}
\newcommand{\pset}{\wp}
\newcommand{\comp}{^{\mathsf{c}}}
\newcommand{\dom}{\operatorname{dom}}
\newcommand{\ran}{\operatorname{ran}}
\newcommand{\id}{\operatorname{id}}
\newcommand{\card}{\operatorname{card}}
\newcommand{\Span}{\operatorname{span}}
\newcommand{\cl}{\operatorname{cl}}
\newcommand{\inte}{\operatorname{int}}
\newcommand{\bd}{\operatorname{bd}}
\newcommand{\lub}{\operatorname{lub}}
\newcommand{\glb}{\operatorname{glb}}
\newcommand{\sy}{\mathbin{\triangle}}
\newcommand{\lequiv}{\sim}            % logical equivalence
\newcommand{\iffdef}{:\Longleftrightarrow}
\newcommand{\T}{\mathrm{T}}
\newcommand{\F}{\mathrm{F}}
\newcommand{\eps}{\varepsilon}
\newcommand{\abs}[1]{\left|#1\right|}
\newcommand{\pair}[2]{(#1,#2)}
\newcommand{\set}[2]{\{\,#1 \mid #2\,\}}
\newcommand{\rest}[2]{#1\!\restriction_{#2}}
\newcommand{\bigcupd}{\bigcup}
